% SECTION_GHOSTS: extension of finite-prefix control to the small primes.
\subsection{A closed local remainder at every coefficient prime}

We use the original A.44 family: $n_i=(i+4)!$, and the supports of
$A=3\prod_iq_i^{n_i}$ and $B=-2\prod_ir_i^{n_i}$ are disjoint.
All continuations in the following statements are required to be legal
and ordinary. The statements control their already existing global
integral class-two sections; they do not construct an all-depth section.

\begin{result}{B.20}{Proved}{Finite-prefix control includes two and three}
\label{coh:smallprefix}
Fix a coefficient prime $\ell$ and an integer $N\geq1$. Put
$M=N+1$ if $\ell=2$, and $M=N$ otherwise. Let $J_M$ be the least
index with $v_\ell(n_{J_M})\geq M$, and define
\[
 S_{\ell,M}=(\{2,3\}\cup\{q_i,r_i:i<J_M\})\setminus\{\ell\},
 \qquad
 D=M+\max_{p\in S_{\ell,M}}v_\ell(p-1).
\]
The maximum over the empty set is zero. Choose $L\geq J_M$ with
$v_\ell(n_L)\geq D$, and fix the prime prefix through $L-1$.
Then every ordinary continuation has the same set of geometric-conjugacy
classes of reductions of its integral pro-$\ell$ class-two sections
modulo $\ell^N$, after using the common abelian cocycle representatives.
At $\ell\geq5$ this set is a singleton. At two and three the assertion
concerns the entire finite set, and does not attach consistent labels
to its individual members. This confirms and extends
GT\_TOWER\_CHECK C.1.21.
\end{result}
\begin{proof}
Fix literal Kummer representatives for the common prefix products.
For each continuation choose the representatives obtained by adding
the uniformly convergent tail sums $\sum_{i\geq L}n_i\kappa_\ell(q_i)$
and $\sum_{i\geq L}n_i\kappa_\ell(r_i)$. Two continuations then
have literal abelian cocycles differing by $\ell^D$ times integral
cocycles. Normalize their class-two sections to those cocycles and
denote their central coordinates by $h_1,h_2$.

At an old $A$-support prime $p\in S_{\ell,M}$, ordinarity gives
$b_\nu=\delta u_\nu$ locally and the exact integral tame relation
gives $h_\nu(T_p)=m_pu_\nu$. On Frobenius,
$(p-1)(u_1-u_2)=(b_1-b_2)(F_p)$ is divisible by $\ell^D$.
Hence $u_1-u_2$ is divisible by $\ell^M$, and the central inertia
coordinates agree modulo $\ell^M$. At an old $B$-support prime use
$a_\nu=\delta v_\nu$ and $h_\nu(T_p)=-m_pv_\nu$; the same
division bound applies. At every other support prime different from
$\ell$, the valuation is divisible by $\ell^M$, so both central
inertia coordinates vanish modulo $\ell^M$. At a zero-valuation
prime the exact tame relation gives $(p^2-p)h_\nu(T_p)=0$, hence
zero integrally. Thus
\[
 w=\overline h_1-\overline h_2
       \in Z^1(G_{\mathbb Q,S},\mathbb Z/\ell^M(2)).
\]
It is a cocycle because the abelian cocycles agree modulo $\ell^M$,
and it descends to $G_{\mathbb Q,S}$ because it kills every inertia
group away from $\ell$.

The coefficient exact sequence gives
\[
 0\to H^1(G_{\mathbb Q,S},T)/\ell^m
 \to H^1(G_{\mathbb Q,S},T_m)
 \xrightarrow{\partial_m}H^2(G_{\mathbb Q,S},T)[\ell^m]\to0,
 \qquad T=\mathbb Z_\ell(2).
\]
The integral first group equals the absolute first group of A.26,
because all its integral cocycles are unramified away from $\ell$.
The integral second group is zero for odd $\ell$ and is
$\mathbb Z/2$ for $\ell=2$, as checked in A.25.
Consequently for odd $\ell$ every finite class $[w]$ is the reduction
of a multiple of $\eta$; for $\ell\geq5$ it is a coboundary.

At two, compare the coefficient exact sequences at exponents $M=N+1$
and $N$. The coefficient reduction is induced by multiplication by two
on the left copy of $T$, and the identity on the middle copy.
Naturality gives
$\partial_N(\overline w)=2\partial_M(w)=0$, because the target
integral second group is killed by two. Thus the class of $w$ modulo
$2^N$ is again the reduction of an integral multiple of $\eta$.
This one extra exponent is why the odd-prime argument cannot simply
be repeated unchanged at two.

Subtract the corresponding integral central twist from the second
section. The remaining central difference modulo $\ell^N$ is an
actual coboundary, hence is central geometric conjugation. Every
integral multiple of $\eta$ is available by A.27. This proves both
inclusions between the finite sets of reductions of the two
continuations, and therefore their equality. The argument neither
chooses nor determines a unique global integral lift above the prefix.
\end{proof}

\begin{result}{B.21}{Proved}{Closedness of the correctly quantified local remainder}
\label{coh:closedlocal}
Let $\mathscr A$ have the legal-prefix topology. For a pair in this
family let $\mathcal S_2(A,B)$ be its finite set of $24$ global integral
class-two section classes. Define
\[
 \mathscr R_{3,\mathrm{loc}}=
 \{(A,B)\in\mathscr A:\exists s\in\mathcal S_2(A,B),\quad
       o_{3,\ell}(s)|_{G_{\mathbb Q_v}}=0
       \text{ for every }\ell\text{ and every }v\}.
\]
This set is closed. Every output outside it has a finite prefix such
that all ordinary legal continuations also lie outside it. The
existential choice of $s$ is made globally, before all local tests;
allowing a separate class-two section at each place gives a different
condition. This result does not identify $\mathscr R_{3,\mathrm{loc}}$
with the global class-three survivor set.
\end{result}
\begin{proof}
Fix an output outside the set. Each of its $24$ classes has a nonzero
local third obstruction at some coefficient prime and some place.
Choose one such witness for each class. Local continuous cohomology
has zero finite-coefficient derived limit, since its finite first
cohomology groups are finite. Thus each nonzero local obstruction is
detected at a finite coefficient exponent. There are finitely many
chosen witnesses and exponents.

We explain why these finite tests depend on a sufficiently long
prefix, including at two. In fixed Hall coordinates the actual
class-three defect is a polynomial in the class-two coordinates,
with bounded rational denominators, and fixed continuous coefficients
from the arithmetic action. For the tripod one may use the explicit
formulas in GT\_TOWER\_CHECK C.1.15, independently checked here;
in the varying section coordinates the only denominator is two.
Consequently agreement modulo $\ell^{n+v_\ell(2)}$ of literal
class-two coordinates implies agreement modulo $\ell^n$ of defects.
A geometric conjugator at a finite class-two level has an integral
lift to class two and then a set lift to class three. Conjugating
before comparing changes the central obstruction by a coboundary,
so the finite cohomology class of the obstruction is preserved.
Applying B.20 at the increased exponent therefore preserves all the
chosen finite local tests as a set of tests on the finite family of
global class-two classes. The extra coefficient step in B.20 takes
care of the two-primary Bockstein.

More explicitly, there is no cross-prime consistency issue: by A.27
the global family is the product of its two-primary and three-primary
families, and is a singleton at every other coefficient prime.
If a global class passed every local test, its components would pass
all the chosen tests at each coefficient prime. Conversely such
component choices combine to a global class by the Chinese remainder
description of $\mathbb Z/24$. Since the original finite collection
excludes the entire global family, and B.20 preserves its component
sets and every test, it excludes the entire family for every ordinary
continuation of a sufficiently long prefix. The complement is open.

The last sentence follows from A.22: global class-three lifting
requires literal membership in $d_3R_\ell$, whereas local
vanishing supplies only coordinatewise divisibility. The additional
condition is that the coordinatewise primitive tends to zero.
No assertion made here evaluates that condition for an arbitrary
member of the closed local remainder.
\end{proof}

\subsection{The actual integral tame morphism}

\begin{result}{B.22}{Proved}{A coefficient-compatible tame cochain representative}
\label{coh:tameedge}
Let $r\geq2$ and let $C_t^\bullet$ be the complex of normalized
continuous $\mathbb Z_\ell(r)$-valued global cochains whose literal
local restrictions at every $q\ne\ell$ factor through
$G_q^{t,\ell}=\mathbb Z_\ell(1)\rtimes\widehat{\mathbb Z}$.
Choose a tame generator $T_q$ and arithmetic Frobenius $F_q$ with
$F_qT_qF_q^{-1}=T_q^q$. The restriction of the morphism
$\Psi_{\ell,r}$ in A.22 to $H^2(C_t^\bullet)$ is induced by
\[
 b\longmapsto(b(T_q))_q,\qquad
 c\longmapsto(z_q(c))_q,
\]
where
\[
 z_q(c)=q^{-1}\left(c(F_q,T_q)-c(T_q^q,F_q)
                       -\sum_{j=1}^{q-1}c(T_q^j,T_q)\right).
\]
The target complex is $[R_\ell\xrightarrow{d_r}R_\ell][-1]$.
In particular this is a formula for the actual continuous residue
class, including its derived-limit contribution. This confirms
GT\_TOWER\_CHECK C.1.14.
\end{result}
\begin{proof}
For a normalized one-cochain $b$, expanding
$\delta b(g,h)=g b(h)-b(gh)+b(g)$ and using $FT=T^qF$ gives
\[
 \sum_{j=1}^{q-1}\delta b(T^j,T)=qb(T)-b(T^q),
 \qquad z_q(\delta b)=(q^{r-1}-1)b(T).
\]
Also $(\delta u)(T)=0$ for a zero-cochain $u$. Extending the
displayed maps by zero in the other degrees therefore gives a
cochain map. It vanishes on normalized unramified cochains.

We identify the resulting map, rather than infer it from the group
isomorphism of A.22. At a finite coefficient exponent, represent a
normalized cocycle $c$ by its extension with multiplication
$(u,g)(v,h)=(u+gv+c(g,h),gh)$. Put $\tau=(0,T)$ and
$\phi=(0,F)$. The central discrepancy in
$\phi\tau=\tau^q\phi$ is precisely $qz_q(c)$, because the
central coordinate of $\tau^q$ is
$\sum_{j=1}^{q-1}c(T^j,T)$. Powers of $\tau$ give a continuous
splitting over the pro-$\ell$ inertia group; its projection is the
identity on $\mathbb Z_\ell$. Conjugating that splitting by
Frobenius and evaluating on $T$ changes the argument to $T^{1/q}$.
Thus the discrepancy of splittings is $q^{-1}$ times the displayed
relator discrepancy. The induced Frobenius action on an inertia
homomorphism $u$ is
$(F\cdot u)(T)=F u(F^{-1}TF)=q^{r-1}u(T)$.
Consequently the two successive cyclic resolutions, first for
inertia and then for Frobenius, give exactly the complex
$[T_n(-1)\xrightarrow{q^{r-1}-1}T_n(-1)][-1]$ with these maps.
In degree one its map is actual restriction to inertia, fixing the
orientation of the Tate coordinate.

The purity description of the relative complex of unramified inside
local cochains is the same two-term residue-field complex used in
A.22. The construction just given is its explicit edge morphism.
The discarded inertia kernel is pro-prime-to-$\ell$; averaging on
its finite quotients is exact on these coefficients. Thus inflation
from the tame quotient gives the same relative morphism for the
local Galois group. Every construction commutes with coefficient
reduction, and with finite localization of the global cochain
complex. Therefore its derived inverse limit is the actual supports
morphism defining $\Psi_{\ell,r}$.

Finally all the displayed families belong to $R_\ell$. At each
finite coefficient exponent a global continuous cochain factors
through a finite Galois quotient. Outside its finite ramification
set, $T_q$ and all $T_q^j$ have trivial image, and normalization
makes every term zero. Division by $q$ is by an $\ell$-adic unit.
The same argument applies to $b(T_q)$. A global coboundary thus
changes the residue by $d_rR_\ell$, not by its closure. This
proves the statement about the derived-limit contribution.
\end{proof}

\begin{result}{B.23}{Proved}{The continuous class-three criterion for a specified lift}
\label{coh:third}
Write a specified global integral class-two section as
$s(g)=y^{b(g)}x^{a(g)}C^{h(g)}s_0(g)$, with $C=[x,y]$ and
$\delta h=-a\cup b$. Set $D=[C,x]$, $F_3=[C,y]$,
$\epsilon=(\chi-1)/2$, and $\eta=(\chi^2-1)/24$.
Its two actual class-three obstruction classes are represented by
\[
 W_D=-(a+\epsilon)\cup h-\binom a2\cup b,\qquad
 W_F=(b+\epsilon)\cup(ab-h)+\binom b2\cup a-\eta\cup b.
\]
Here products and binomial expressions denote continuous cochains
with the indicated weight, rather than assertions that they are
one-cocycles. At $q\ne\ell$ put
\[
 \begin{gathered}
 (m,n,t)=(a(T_q),b(T_q),h(T_q)),\quad
 (\alpha,\beta,\gamma)=(a(F_q),b(F_q),h(F_q)),\\
 e=(q-1)/2,\qquad e_2=(q^2-1)/24,
 \end{gathered}
\]
and set
$B_q(z)=q^{-1}\sum_{j=1}^{q-1}\binom{jz}{2}
       =(q-1)((2q-1)z^2-3z)/12$.
Their images under $\Psi_{\ell,3}$ are the classes of the actual
$R_\ell$-sequences
\begin{align*}
 Z_{D,q}={}&m\gamma-q(\alpha+e)t-n\binom\alpha2
       +q^{-1}\beta\binom{qm}{2}+emt+nB_q(m),\\
 Z_{F,q}={}&n(\gamma-\alpha\beta)+q(\beta+e)(mn-t)
       +m\binom\beta2-e_2n-q^{-1}\alpha\binom{qn}{2}\\
       &{}-e(mn^2-nt)-mB_q(n).
\end{align*}
The section lifts to the integral class-three quotient if and only if,
for every $\ell$, both sequences $Z_{D,q}/(q^2-1)$ and
$Z_{F,q}/(q^2-1)$ are integral and tend to zero $\ell$-adically.
This confirms GT\_TOWER\_CHECK C.1.15.
\end{result}
\begin{proof}
The actual path action has weight-two path coefficient $\eta$, in the
normalization of \cite[(12.12), (12.19), (12.35)]{CohWickelgren2012}.
The published formulas use its letters $(b,a)$ for our $(a,b)$.
Collection with $[u,v]=uvu^{-1}v^{-1}$ gives
\[
 x^a y^b=y^b x^a C^{ab}
       D^{b\binom{a+1}{2}}F_3^{a\binom{b+1}{2}}.
\]
Thus, on a zero-last-coordinate lift, the actual arithmetic action is
\[
 g(y^b x^a C^h)=y^{\chi b}x^{\chi a}C^{\chi^2h}
           D^{-\epsilon\chi^2h}
           F_3^{-\epsilon\chi^2h-\eta(g)\chi b}.
\]
These are integral identities at two and three as well; $\chi^2-1$
is divisible by $24$ in $\widehat{\mathbb Z}$. Substituting the
class-two cocycle equation in the zero-last-coordinate defect gives
\begin{align*}
 c_D(g,k)={}&h(g)\chi(g)a(k)
    +\binom{a(g)+1}{2}\chi(g)b(k)
    +a(g)\chi(g)^2a(k)b(k)-\epsilon(g)\chi(g)^2h(k),\\
 c_F(g,k)={}&h(g)\chi(g)b(k)
    +a(g)\binom{\chi(g)b(k)+1}{2}
    -\epsilon(g)\chi(g)^2h(k)-\eta(g)\chi(g)b(k).
\end{align*}
Direct expansion gives the integral coboundary comparisons
$W_D=c_D+\delta(ha)$ and
$W_F=c_F+\delta(hb-a\binom b2)$. These also agree with
\cite[Propositions 13 and 17, author pp.~17, 19]{CohWickelgren2012} after the stated
change of names. Hence they represent the actual extension defect,
and are two-cocycles.

At $q\ne\ell$ the displayed action is unramified. The section
kills wild and prime-to-$\ell$ tame inertia: on that kernel its
geometric cocycle is a homomorphism from a pro-prime-to-$\ell$ group
to a pro-$\ell$ group. Therefore these literal cochains lie in
the subcomplex of B.22. Substitution of
$a(T^j)=jm$, $b(T^j)=jn$, $\chi(T^j)=1$ into its tame formula
gives the two displayed residues. The sums of the linear and
quadratic powers of $j$ give $B_q$. Its sum expression proves
$\ell$-integrality for every $q\ne\ell$, even where its rational
polynomial expression displays denominators divisible by $\ell$.
The independent symbolic checks are retained in
\path{code/ghost_audit_checks.py}; that check also handles $q=2,3$.

B.22 shows that the residue families are in $R_\ell$ and represent
the actual continuous classes. Since $\Psi_{\ell,3}$ is an
isomorphism by A.22, vanishing is exactly membership in $d_3R_\ell$.
Each nonzero integer $q^2-1$ acts injectively on $\mathbb Z_\ell$,
so this is precisely integrality and convergence of the two unique
coordinatewise quotients. Continuous primitives then correct the
two central coordinates and yield a genuine lift, by A.28. Products
over $\ell$ give the integral assertion.
\end{proof}

\subsection{The stronger boundary-matching gate}

\begin{result}{B.24}{Proved}{Boundary matching forces the global third convergence condition}
\label{coh:boundarygate}
For an ordinary A.44 output, let $s_*$ be the unique global integral
class-two section class which agrees with its actual local point
section in the pro-two quotient at $q=3$ and in the pro-three quotient
at $q=2$. This class exists. Every full pronilpotent section above
the pair has class-two image $s_*$.

Fix $\ell$ and normalize $s_*$ to literal abelian cocycles $a,b$.
At an $A$-support prime $q\ne\ell$, write
$b=\delta u_q$, so $\beta=(q-1)u_q$, and put
\[
 R_{A,q}=\frac{\gamma-q\alpha u_q}{q^2-1}.
\]
At a $B$-support prime $q\ne\ell$, write
$a=\delta v_q$, so $\alpha=(q-1)v_q$, and put
\[
 R_{B,q}=\frac{\gamma+\beta v_q-\eta(F_q)}{q^2-1}
       =\frac{\gamma+\beta v_q}{q^2-1}-\frac1{24}.
\]
Then $s_*|_{G_q}$ is the actual local point class if and only if
the corresponding $R_{A,q}$ or $R_{B,q}$ is in $\mathbb Z_\ell$.
These conditions, at every support $q\ne\ell$ and every $\ell$,
are necessary for a full pronilpotent survivor. If they all hold,
$s_*$ lifts to the global integral class-three quotient: the strict
convergence condition in B.23 follows automatically. This is a
necessary all-depth gate and a sufficient class-three test; it is
not asserted sufficient at every depth.
\end{result}
\begin{proof}
By A.27 choose any global class-two class. Its difference from the
prescribed local point class lies in local weight-two first cohomology.
The two restriction maps in A.29 identify the global two-primary and
three-primary factors with the local groups at $(q,\ell)=(3,2)$ and
$(2,3)$, respectively. Thus exactly one global twist makes both
differences zero. A.31, applied at these two nonzero valuations,
forces the class-two image of any full pronilpotent survivor to be
this class.

We identify the local point class in the actual action. At an
$A$-support prime the inertia image of a point section is conjugate
to $x^m$, where $m=v_q(A)>0$, by the proof of A.31. After that
conjugation its Frobenius coordinate centralizes $x^{qm}$, and is
therefore a power of $x$. Conjugating in class two by $y^{u_q}$,
using the convention $c'(g)=k^{-1}c(g)g(k)$, changes its
coordinates to
\[
       (a(g),\, (\chi(g)-1)u_q,\, \chi(g)a(g)u_q).
\]
Thus the normalized local point has central Frobenius coordinate
$q\alpha u_q$.

At a $B$-support prime its inertia is conjugate to $y^n$.
Here the actual action is
$g(y)=f_g^{-1}y^{\chi(g)}f_g$, with
$f_g\equiv C^{\eta(g)}$ in class two. The tame relation says that
the Frobenius coordinate times $f_{F_q}^{-1}$ centralizes $y^{qn}$.
Consequently the normalized boundary point section before restoring
the first coordinate has central Frobenius coordinate $\eta(F_q)$.
Conjugation by $x^{v_q}$ restores $a=\delta v_q$ and subtracts
$\beta v_q$ from that coordinate. Thus the local point has
$\gamma=\eta(F_q)-\beta v_q$, modulo central coboundaries.
The path term is essential at two and three.

The literal central inertia coordinates of every integral class-two
section over these abelian cocycles are, respectively, $m u_q$ and
$-n v_q$, by the tame relation. Its difference from the normalized
point section is therefore unramified. Central coboundaries change
the Frobenius value by $(q^2-1)\mathbb Z_\ell$. Also local
$H^0(G_q,\mathbb Z_\ell(1)^2)=0$, so a geometric conjugation
between normalized sections is central. These facts prove the two
exact matching criteria. A.31 proves their necessity at every support
prime for an all-depth survivor.

For sufficiency at class three, substitute these local coordinates
in B.23. At an $A$-support prime the exact expressions are
\[
 \frac{Z_{D,q}}{q^2-1}=m R_{A,q}+u_q\binom m2,
 \qquad
 \frac{Z_{F,q}}{q^2-1}=-m\binom{u_q+1}{2}.
\]
At a $B$-support prime they are
\[
 \frac{Z_{D,q}}{q^2-1}=n\binom{v_q+1}{2},
 \qquad
 \frac{Z_{F,q}}{q^2-1}=n R_{B,q}-v_q\binom n2.
\]
The $-n/24$ of B.23 is exactly absorbed in the actual point-class
shift in $R_{B,q}$. These identities hold also at $q=2,3$, with
the equal coefficient prime excluded as stated. At zero-support
primes, $m=n=0$ and the integral tame relation gives $h(T_q)=0$,
so both residues are zero.

Under the matching conditions all $u_q,v_q,R_{A,q},R_{B,q}$ are
integral. The valuations $m,n$ tend to zero $\ell$-adically;
in the A.44 family this also follows directly from the factorial
exponents. The displayed binomials are integral, and
$v_\ell(\binom m2)\geq v_\ell(m)-v_\ell(2)$, with the same
bound for $n$. Hence both normalized residue sequences tend to
zero. B.23 gives an actual global integral class-three lift. No
choice of lifts at higher classes has been made in this argument.
\end{proof}

\begin{result}{B.25}{Proved}{Boundary matching leaves a fixed-ramification obstruction}
\label{coh:higherboundarygate}
Let $r\geq2$, fix $\ell$, and let $s_r$ be a continuous section of
the actual pro-$\ell$ class-$r$ quotient with an ordinary completed
abelian pair. Assume that, at every support prime $q\ne\ell$,
its local class equals the class-$r$ image of the corresponding
ordinary point section. Then its next obstruction has the form
\[
 o_{r+1}(s_r)=\operatorname{inf}(\omega_{r+1}),\qquad
 \omega_{r+1}\in
 H^2(G_{\mathbb Q,\{\ell,\infty\}},
          \mathbb Z_\ell(r+1))^{b_{r+1}},
\]
with a unique $\omega_{r+1}$. In particular it lifts if and only
if this fixed-ramification class is zero. The latter group is finite.
For $r=2$ it is zero, so the local boundary-matching hypothesis
alone implies a global class-three lift for every ordinary pair,
including negative-valuation support patterns.

This removes the strict tame residue contribution under the stated
matching hypothesis. It does not remove the entire absolute
derived-limit kernel: a fixed-ramification class can inflate to a
nonzero class invisible at every finite absolute coefficient level,
as in A.23.
\end{result}
\begin{proof}
Use the actual action and the continuous Hall set section from class
$r$ to class $r+1$, with new central coordinates zero. Its defect
is a vector of continuous two-cocycles $c$ representing the
obstruction. At every $q\ne\ell$ the action is unramified by
the proof of A.31, and the section kills the prime-to-$\ell$
inertia kernel. Thus B.22 applies to each coordinate of $c$.

At a support prime the assumed geometric conjugacy in class $r$
has a lift to the free pro-$\ell$ group. Conjugate the actual full
local point section by this lifted element. The resulting local
section has literal class-$r$ image $s_r|_{G_q}$. Its difference
from the zero-last-coordinate set lift is a central cochain. With
the appropriate sign write it as $f_q$, so $\delta f_q=c|_{G_q}$.
On the tame generator, $f_q(T_q)$ is the negative of the new Hall
coordinate of that actual local section's inertia image.

By the proof of A.31 this inertia image is conjugate in $F_\ell$
to one of
\[
 x^{m_q},\qquad y^{n_q},\qquad ((xy)^{-1})^{-m_q},
\]
in the three valuation patterns $(m_q,0)$, $(0,n_q)$, and
$(m_q,m_q)$ with $m_q<0$, respectively. There is a uniform
bound on all Hall coordinates through class $r+1$ of such
conjugates. Indeed multiplication, inversion, and powering in the
fixed Hall coordinates are rational polynomials
\cite[\S\S3--4]{CohCantEick2019}. The polynomial for a coordinate
of $k^{-1}z^e k$, for any one of the three fixed meridians $z$,
vanishes identically at $e=0$. It is therefore $e$ times a
rational polynomial in $e$ and the integral coordinates of $k$.
After clearing the finitely many denominators there is an integer
$D_r\geq1$, independent of $q,e,k,\ell$, such that
\[
       f_q(T_q)\in e_qD_r^{-1}\mathbb Z_\ell^{b_{r+1}},
\]
where $e_q$ is the corresponding ordinary integer exponent.
The actual values are integral as well. Because $e_q\to0$
$\ell$-adically, these local primitive inertia values tend to zero.
In particular no division by $q^i-1$ occurs in this estimate.

At a zero-support prime the inertia image of $s_r$ is the identity.
To see this, its weight-one coordinate is zero. If its first
possibly nonzero central coordinate has weight $i\geq2$, the
tame relation and the graded action give
$(q^i-q)V_i=0$. Torsion freeness kills $V_i$; induction proves
the assertion. Hence the section is unramified. Choose any lift of
its arithmetic Frobenius element to class $r+1$. A chosen element
in a profinite group defines a continuous homomorphism from
$\widehat{\mathbb Z}$ by profinite powering, and its projection
here is the identity on $\widehat{\mathbb Z}$. This gives an
unramified lifted section, so we choose $f_q(T_q)=0$.

The entire family $(f_q(T_q))_{q\ne\ell}$ consequently belongs
to $R_\ell^{b_{r+1}}$. The cochain identity of B.22 gives
\[
 z_q(c)=(q^r-1)f_q(T_q).
\]
It follows that $\Psi_{\ell,r+1}(o_{r+1}(s_r))=0$ in the
strict quotient $R_\ell/d_{r+1}R_\ell$, with actual image and
without taking a closure. Exactness and injectivity in A.22 give
the unique class $\omega_{r+1}$ and the stated criterion.

For clarity, the finiteness assertion includes the precise real
two-primary alternatives. Put $w=r+1$ and $K_w=K_{2w-2}(\mathbb Z)$,
which is finite. For odd $\ell$, \cite[VI.8.2]{CohWeibelKbook}
identifies the fixed-ramification group with $K_w\{\ell\}$.
For $\ell=2$, \cite[VI.9.11, author p.57]{CohWeibelKbook}
gives the same identification when $w\equiv1,2,3\pmod4$:
in its $w\equiv3$ case, $0\leq\rho<r_1=1$ forces $\rho=0$.
When $w\equiv0\pmod4$, that theorem identifies $K_w\{2\}$
with the kernel of restriction to
$H^2(G_{\mathbb R},\mathbb Z_2(w))=\mathbb Z/2$.
Thus in this case there is an exact sequence
\[
 0\to K_w\{2\}\to H^2(G_{\mathbb Q,\{2,\infty\}},
                  \mathbb Z_2(w))\to J_w\to0,
       \qquad J_w\subseteq\mathbb Z/2,
\]
which is sufficient for finiteness without any extra surjectivity claim.
These comparisons apply to the ring with $\ell$ inverted:
localizing $\mathbb Z$ at $\ell$ does not change its
$\ell$-primary group, since the finite-field residue in degree
$2w-3$ has order $\ell^{w-1}-1$, prime to $\ell$
\cite[IV.1.13; V.6.8]{CohWeibelKbook}.
The weight-three group vanishes by the fully two-primary-inclusive
calculation in A.22. Finally A.23 demonstrates why the last
qualification in the statement is necessary.
\end{proof}

\begin{result}{B.26}{Proved}{The separate obstruction to preserving boundary matching}
\label{coh:matchingobstruction}
Retain the hypotheses and notation of B.25, and suppose
$\omega_{r+1}=0$. Put $T=\mathbb Z_\ell(r+1)^{b_{r+1}}$ and
let $\Sigma$ be the support primes different from $\ell$.
Choose one global class-$(r+1)$ lift $t$ of $s_r$. There is a
well-defined tuple of local discrepancy classes
\[
 d(t)=(d_q(t))_{q\in\Sigma}\in
       \prod_{q\in\Sigma}H^1(G_q,T)
\]
such that a global lift matching the actual point class at every
$q\in\Sigma$ exists if and only if
\[
 d(t)\in\operatorname{im}\left(
 H^1(G_{\mathbb Q},T)\longrightarrow
                     \prod_{q\in\Sigma}H^1(G_q,T)\right).
\]
The image is compact and closed. If this condition fails, some finite
set of support primes and some finite coefficient exponents already
exclude every integral global twist of $t$. If it holds, the matching
lifts form a torsor under the kernel of the displayed localization map.
This is an additional condition after the finite second-cohomology
obstruction in B.25 has vanished.
\end{result}
\begin{proof}
We first remove the stabilizer ambiguity. For any actual class-$r$
section $c$ over $G_{\mathbb Q}$, an element $k$ of $N_{r,\ell}$
centralizing its image is fixed under the twisted action
$k\mapsto c(g)g(k)c(g)^{-1}$. If $k$ has a first nonzero central
coordinate in weight $i$, inner conjugation has no effect on this
coordinate, and invariance gives
$(\chi(g)^i-1)k_i=0$. There is a $g$ with $\chi(g)^i\ne1$,
and the coordinate module is torsion free, a contradiction. The same
argument over $G_q$, $q\ne\ell$, uses arithmetic Frobenius and
the nonzero integer $q^i-1$. Thus both global and these local
geometric stabilizers are trivial.

Fix a literal representative of $s_r$. At $q\in\Sigma$, choose
an actual full point section representing the prescribed local point.
By hypothesis its class-$r$ image is geometrically conjugate to
$s_r|_{G_q}$, and the conjugating element in $N_{r,\ell}$ is
unique by the stabilizer calculation. Lift that element to
$N_{r+1,\ell}$ and conjugate the local point section. Its
class-$r$ image now agrees literally with $s_r|_{G_q}$.
Two choices of the lifted conjugator differ by an element of the
central kernel $T$, so the resulting local lifted sections differ
by a central coboundary. Changing the initial representative of the
point section has the same effect after alignment. Consequently its
class in the central lift fibre is unambiguous.

Two lifts with the same literal class-$r$ image differ by a
continuous $T$-valued one-cocycle, since the inner action of class
$r$ is trivial on $T$. Modulo central conjugation their difference
is a class in $H^1(G_q,T)$. Define $d_q(t)$ as the class which
must be added to $t|_{G_q}$ to obtain the aligned point lift.
The same argument globally, together with the vanishing of the
global stabilizer, identifies all global lift classes above $s_r$
with the torsor $t+H^1(G_{\mathbb Q},T)$. Its localization sends
$t+h$ to the desired point class precisely when
$\operatorname{res}_q(h)=d_q(t)$ for every $q$. This proves the
criterion and the assertion about the torsor of matching lifts.
Replacing $t$ changes $d(t)$ by an element of the displayed image,
so solvability is independent of that choice.

For completeness, the compactness here has a fixed coefficient
module and does not permit fresh ramification in finite reductions.
Integral weight-$(r+1)$ cocycles are unramified outside $\ell$,
by the tame relation, as proved in A.20. Their compact-open cocycle
module is compact. The coboundaries are the continuous image of
$T$, a compact module, and hence form a closed subgroup. Therefore
$H^1(G_{\mathbb Q},T)$ is compact in the quotient topology.
At $q\ne\ell$ the same tame relation kills inertia for integral
cocycles, giving
\[
 H^1(G_q,T)=
 \left(\mathbb Z_\ell/(q^{r+1}-1)\mathbb Z_\ell\right)^{b_{r+1}}.
\]
These groups are finite, and localization is continuous: in unramified
coordinates it is evaluation at Frobenius followed by a finite
quotient. Its image in their product is thus compact and closed.

If the discrepancy tuple is outside that image, a basic open
neighbourhood of it disjoint from the image involves only finitely
many coordinates. These give the finite set of places. At such a
place, choose a coefficient exponent at least
$v_\ell(q^{r+1}-1)$. The coefficient exact sequence shows that
reduction of the displayed integral first-cohomology group is
injective: its kernel is its multiple by that power of $\ell$,
which is zero. The same finite separation is therefore visible in
finite coefficient cohomology. This last assertion concerns the
reductions of integral global twists. Replacing those twists by all
finite-coefficient global cocycles would change the problem, since
those cocycles may introduce new ramification.
\end{proof}

The exact tower problem can consequently be expressed using B.25
and B.26 together. At each class choose a global section which
matches the boundary point classes, first annul its finite
fixed-ramification second-cohomology class, and then test the
localization condition for the differences of its lifts. Neither
step implies the other. For a fixed ordinary pair the sets of such
sections are compact by A.20--A.21. If they are nonempty at every
class, the compact finite-intersection argument supplies one
compatible sequence and a full pro-$\ell$ section. Conversely
A.31 forces every full pro-$\ell$ survivor into these sets. Taking
all coefficient primes gives the same exact criterion for the full
pronilpotent quotient. The unresolved question is whether a ghost
can satisfy these successive global and matching conditions at all
classes simultaneously.
